\section{Composition: the literal 228-state table has the required calibration}
\label{sec:composition}

Let $T$ be the released K230 source tree and let $T'$ be the supplied
5,121-leaf tree.  Let $M$ be the 228-state literal table, first considered
with its retained K230 state labels.  Its initial state is 0, its dispatcher
is 84, and its separate halting outcome has retained label 230.
The final contiguous renaming is a bijection on working states and maps
the separate halting outcome to the separate halt label; it therefore
preserves each transition and blank-tape halting by induction on time.
The argument below concerns the complete table, including its shared
initialization and decoder states.

\subsection{The source transformation preserves the whole program}

By Proposition~\ref{prop:source}, the four closed substitutions preserve
source halting from every register file. For completeness, the lexical
control-scope calculation connecting that macro lemma to the complete
structured program is as follows.

Both fragments have boundary signature $(B,C,H)=(0,0,\mathrm{false})$:
no break or continue escapes, and neither can halt.  For a binary node
with child signatures $(B_0,C_0,H_0)$ and $(B_1,C_1,H_1)$, the signature is
\[
 \left(\max(B_0\mathbin{\dot-}1,B_1),
       \max(C_0,C_1\mathbin{\dot-}1),H_0\lor H_1\right),
\]
where $n\mathbin{\dot-}1=\max(n-1,0)$.
This calculation checks both entire source trees to have signature
$(0,0,\mathrm{true})$.  Their full binary prefix listings also prove
finite decoding, and all control depths are positive.  Thus the generic
well-formed source semantics applies to both programs.

Replacing a terminating closed invocation inside a finite structured run
preserves its exit mode and entire returned register file.  Induction on
the surrounding finite run derivation proves this in either direction.
The same replacement preserves every completed normal execution of the
main tree, so it also preserves any finite sequence of main-tree
iterations preceding a halt.  The inherited structured/path equivalence
\texttt{rm\_halts\_iff\_iter} consequently gives
\begin{equation}
 \operatorname{RMHalt}(T,[],R)
       \longleftrightarrow\operatorname{RMHalt}(T',[],R)
 \quad\text{for every register file }R.
 \label{eq:k228-source-equivalence}
\end{equation}
Here $\operatorname{RMHalt}$ means reaching the explicit halt leaf;
it does not identify a malformed or stuck source configuration with a
halt.\footnote{Inherited generic definitions and equivalence:
\url{https://github.com/elliotglazer/k230-rocq/blob/3142e4dd980f0be12f4abb636c4ae2528afc17eb/coq/tm/RM.v}
and
\url{https://github.com/elliotglazer/k230-rocq/blob/3142e4dd980f0be12f4abb636c4ae2528afc17eb/coq/tm/RMEquiv.v}.
The release attributes these reused generic semantics to Jeff's reference
development.  The exact local replacement and source identity are
rechecked by \texttt{code/check\_composition.py}.}

\subsection{Why the universal leaf cases can be iterated}

Write $B(v)=0\,1^{v+1}$ and extend $B$ to finite lists by concatenation.
Let $\pi R$ put the nine logical registers in physical order, using
\[
        \pi=(6,2,0,4,3,7,8,5,1)
        \quad\text{(logical index $i$ is stored at physical index $\pi_i$).}
\]
A dispatch frame $F(P,G,R,E)$ has state 84, a single 1 immediately to
the left of the head followed by the all-zero left exterior, and right word
\begin{equation}
        P\,0^{G-|P|}\,1\,B(\pi R\mathbin{+\!+}E)\,0^\omega,
        \qquad G\ge64,\quad E\ne[].
 \label{eq:k228-frame}
\end{equation}
Every entry of $R$ and of the nonempty finite spare list $E$ is an arbitrary
natural number.  The pointer $P$ has implicit trailing zeros.  It can end
at an internal node, or contain trailing zeros beyond a shallower leaf.

The crucial closure fact is the following exact finite-word identity.
Let $L$ be the leaf reached by reading $P0^\omega$.  For every reachable
source pointer, all bits of $P$ beyond $L$ are zero; hence
\begin{equation}
       P\,0^{G-|P|}=L\,0^{G-|L|}.
 \label{eq:k228-padding}
\end{equation}
The inherited generic lemmas \texttt{no\_stray} and \texttt{split\_tape}
prove precisely this property.  It also has a short direct proof.
Stripping bits from a leaf yields a prefix of that leaf, which is a
tree node until decoding reaches a leaf.  Incrementing such a prefix
either gives only zeros on overflow, or changes $A0\,1^m$ to
$A1\,0^m$ at an existing binary branch.  Every subsequent excess bit
is zero.  This covers increment, decrement in either branch, break,
and continue.\footnote{The exact inherited invariant and word-splitting
lemma are in
\url{https://github.com/elliotglazer/k230-rocq/blob/3142e4dd980f0be12f4abb636c4ae2528afc17eb/coq/tm/Sim230.v};
\texttt{no\_stray} is in the preceding \texttt{RMEquiv.v} source.}

For additional exact source binding, the independent composition checker
evaluates the successor of every nonhalting leaf/guard case and checks
\eqref{eq:k228-padding} directly.  All 6,655 identities pass: 4,344
successors extend with zeros to the next leaf, 1,781 already equal that
leaf, and 530 discard only trailing zeros.  The sole remaining case is
the explicit halt leaf.  The empty pointer decodes to leaf
\texttt{000}, so the initialization case has the same property.

The symbolic runtime proof covers all 6,656 leaf/guard cases, with the
decrement guard split into zero and an arbitrary positive value.
Its parameters are nine independent register values, an arbitrary extra
gap $G-64$, an arbitrary first spare, and an arbitrary finite spare tail.
Literal steps use the exact table.  Self-loop scans are justified by
induction on a unary run length; four exact one-block calculations extend
to arbitrary spare lists by list induction.  These rules preserve the
entire tape word and the all-zero exterior.  Their finite derivations
give finite concrete runs for every parameter assignment.  The bound
4,000 limits construction of each symbolic proof, with no bound on
concrete values or runtime.

For a nonhalting leaf, the proved target is
\[
             F(P',G+\delta+o,R',E'),
\]
where $P'$ and $R'$ are the exact source successor, $\delta=1$ only for
a positive decrement, and $o$ is the binary counter overflow bit.
An increment appends a spare of value 1; all other operations preserve
the spare list.  In particular, it stays nonempty.  The proved head
displacement is $-o$, so an overflow moves the left marker as well as
increasing the total instruction area.  All three overflow cases are
included in the independent check.  Because $G+\delta+o\ge64$ and every
new source leaf has depth at most 40, \eqref{eq:k228-padding} converts
this exact target to another admissible leaf frame.  This proves closure
of the runtime invariant, including changes in pointer length and return
counter overflow.

\subsection{Initialization, finite progress, and the halting equivalence}

The complete blank run of $M$ reaches state 84 after exactly 409,835
transitions.  Its full tape support is the frame
\[
 F([],181,R_0,[1]^{167}),\qquad
              R_0=[3,0,0,0,3,0,0,0,0].
\]
Thus the extra gap is 117.  The nine physical values are
$[0,0,0,3,0,0,3,0,0]$, exactly the inherited K230 boot values.
The independently implemented dense and sparse simulations agree on
the full finite support (351 one-cells), head position $-5$, and state;
all other cells are zero.  Absolute translation of the frame has no
semantic effect.  This reaches the actual nonzero seed used by the
inherited enumerator completeness theorem.

Every checked runtime segment has at least 13 concrete transitions.
A nonhalting segment cannot pass through the separate halt outcome:
the literal table has no transition there, so the certificate could not
continue to its nonhalting target.  The one halt-leaf certificate reaches
that outcome in a finite positive number of transitions.  State 0 can
appear inside such segments because it also serves the shared runtime;
its literal transitions are included without treating a revisit as a new
boot or as a halting boundary.  All 456 table rows are covered by the
universal certificates.  This coverage is an audit fact; soundness comes
from the stated induction rules and complete case partition.

Induction over source steps now concatenates the finite runtime
segments.  If the source reaches its halt leaf, the machine halts.
Conversely, if the source never reaches that leaf, full decoding and the
preserved frame provide another finite nonhalting segment at every
boundary.  After $n$ segments the elapsed time is at least $13n$;
the boundaries therefore exceed every finite time.  No finite-time halt
can occur between them, and no one segment can diverge.  This proves
\begin{equation}
 \operatorname{Halt}(M,\mathrm{blank})
        \longleftrightarrow\operatorname{RMHalt}(T',[],R_0).
 \label{eq:k228-runtime-equivalence}
\end{equation}

The same runtime argument applies directly to the original source $T$
and the pinned literal K230 table. The separately completed baseline
certificate covers all 6,660 guarded cases, with the same positive
progress bound and frame semantics. Its exact 408,419-transition boot
reaches $F([],181,R_0,[1]^{167})$. Therefore
\[
 \operatorname{Halt}(K230,\mathrm{blank})
 \longleftrightarrow\operatorname{RMHalt}(T,[],R_0).
\]
This direct check avoids relying on an unverified identification between
the released site JSON and an independently represented Rocq program AST.

Combining \eqref{eq:k228-runtime-equivalence},
\eqref{eq:k228-source-equivalence}, and the inherited K230 calibration
gives the completed ordinary mathematical theorem
\begin{equation}
 \operatorname{Halt}(K228,\mathrm{blank})
       \longleftrightarrow\operatorname{Halt}(K230,\mathrm{blank})
       \longleftrightarrow\texttt{StandardZFC}\vdash\bot.
 \label{eq:k228-complete}
\end{equation}
The inherited proof of completeness from $R_0$ is used as stated; no
all-zero seed or bounded search assumption is substituted for it.

\paragraph{Meaning of ``provable in ZFC.''}
Choose a standard effective coding of finite first-order ZFC derivations,
and define $\operatorname{Con}(\mathrm{ZFC})$ as the absence of a finite
contradiction proof.  The new argument uses finite words, finite tables,
finite lists and natural-number induction, all formalizable in ZFC.
The inherited logical reductions act on finite object-level derivations:
the lower interpretation translates a ZFC contradiction into a K
contradiction, and the upper construction translates a K contradiction
into a finite ZFC derivation using the required finite schema instances.
In particular, the constructible-hierarchy argument is represented by
object-level proof transformations, not by assuming the existence of a
model of ZFC.\footnote{The finite-axiom collection and the actual
derivation transfer can be read directly in
\texttt{KPrimeOrdIR\_standardZFC\_inconsistent}, in
\url{https://github.com/elliotglazer/k230-rocq/blob/3142e4dd980f0be12f4abb636c4ae2528afc17eb/coq/settheory/KUpperConsistency.v}.
The other direction is \texttt{KPrimeOrd\_standardZFC\_inconsistent}, in
\url{https://github.com/elliotglazer/k230-rocq/blob/3142e4dd980f0be12f4abb636c4ae2528afc17eb/coq/settheory/KQuotZFC.v}.}
These particular syntactic constructions and their
correctness arguments are likewise formalizable in ordinary ZFC.
Consequently \eqref{eq:k228-complete}, followed by negation, yields
\[
 \boxed{\quad
 \mathrm{ZFC}\vdash
 \bigl(\operatorname{Nonhalt}(K228,\mathrm{blank})
              \longleftrightarrow\operatorname{Con}(\mathrm{ZFC})\bigr).
 \quad}
\]
This is the scope of the mathematical proof.  It does not assert
$\operatorname{Con}(\mathrm{ZFC})$, require a general soundness theorem
for Rocq inside ZFC, or assert a new machine-checked internal PA proof.
The inherited K230 endpoint is the released Rocq theorem; the new
228-state transfer is supplied here as an ordinary proof with executable
finite and symbolic certificates.  A freshly kernel-checked Rocq theorem
for K228, or a fully encoded object-language ZFC/PA derivation, is a
further verification artifact and is not claimed by these checks.
